\documentclass[11pt]{article}
\usepackage{amsmath,amssymb,amsthm,hyperref}
\newtheorem{theorem}{Theorem}
\newtheorem{proposition}[theorem]{Proposition}
\newtheorem{corollary}[theorem]{Corollary}
\title{Full-Cube Constraints in Lower-Bound Continuations}
\author{}
\date{November 23, 2026}
\begin{document}
\maketitle

\paragraph{Status.}
The statements here are complete written arguments. They are not Lean
declarations and do not establish a new asymptotic bound for unrestricted
\texttt{HeadComplexity.worstCaseHeadComplexity}. The attention model,
head count, and denominator admissibility remain unchanged. The purpose
is to delimit two concrete counting mechanisms before using them in a
claimed lower bound. Neither statement excludes arbitrary target-dependent
denominators, numerator cancellation, or asymmetric reciprocal queries.

\section{A full-cube obstruction to compressing symmetric Cauchy sheets}

The construction in \texttt{compact-cauchy-sheets.tex} gives factorially
many regular zeros of a system of reciprocal equations. It uses one-hot
row coordinates to supply independently chosen affine numerator values.
Replacing those row coordinates by binary addresses does not preserve
that freedom. The following obstruction applies to the most direct
attempt at doing so.

\begin{theorem}[Strict concavity of the reciprocal Cauchy sum]
Let $H\geq 2$, let $w_1,\ldots,w_H>0$, and let
$\lambda_1,\ldots,\lambda_H$ be real numbers with at least two distinct
values. For $d>\max_h\lambda_h$, put
\[
 S(d)=\sum_{h=1}^H\frac{w_h}{d-\lambda_h},\qquad R(d)=\frac1{S(d)}.
\]
Then $R''(d)<0$ throughout its domain.
\end{theorem}

\begin{proof}
Set $T_j(d)=\sum_h w_h/(d-\lambda_h)^j$. Direct differentiation gives
$S'=-T_2$ and $S''=2T_3$, hence
\[
 R''(d)=\frac{2T_2(d)^2-2T_1(d)T_3(d)}{T_1(d)^3}.
\]
Apply Cauchy--Schwarz to the two vectors with coordinates
$\sqrt{w_h/(d-\lambda_h)}$ and
$\sqrt{w_h/(d-\lambda_h)^3}$. Equality would require the ratio
$1/(d-\lambda_h)$ to be independent of $h$. The stated distinctness
excludes equality, so $T_2^2<T_1T_3$. Every denominator is positive,
and the displayed formula proves the assertion.
\end{proof}

\begin{corollary}[No two-bit affine encoding of all symmetric zero rows]
Let $D,A$ be affine functions on a Boolean cube. Suppose that $D$
takes values larger than $\max_h\lambda_h$ throughout that cube.
If two coordinate slopes of $D$ are nonzero and have the same sign,
then the score
\[
 F(u)=A(u)\sum_{h=1}^H\frac{w_h}{D(u)-\lambda_h}-1
\]
cannot vanish on every vertex of any two-dimensional face varying
those coordinates.
\end{corollary}

\begin{proof}
Fix all other coordinates and write the four denominator arguments as
$d,d+a,d+b,d+a+b$, where $ab>0$. If all four scores vanished,
their positive denominators would imply $A(u)=R(D(u))$ at all
four vertices. The mixed finite difference of an affine function
is zero. The corresponding difference of $R(D(u))$ is
\[
 R(d+a+b)-R(d+a)-R(d+b)+R(d)
 =ab\int_0^1\int_0^1 R''(d+sa+tb)\,ds\,dt<0.
\]
The whole rectangle lies in the domain, since the affine argument
has its extreme values at its vertices. This is a contradiction.
\end{proof}

This is compatible with actual attention admissibility: all slopes
can be strictly negative and every denominator can exceed $3/4$ on
the whole cube. The obstruction is the affine relation among row
numerator values, not a pole or a failure of positivity. It does not
apply once head-dependent row numerators or row denominators break
the displayed symmetric form.

\section{Signed residues restore factorially many components}

The previous obstruction does not extend to arbitrary numerators.
The following construction is a new full-cube-compatible counterexample
to a parameter-linear bound on logarithms of intersection-component
counts. It uses cancellation by signed residues.

\begin{theorem}[Binary-addressed interpolation components]
Let $d_1,\ldots,d_H$ be distinct real numbers, all larger than $\beta>0$.
For $b\in(0,\beta)^H$ and $c\in\mathbb R^H$, set
\[
 F_i(b,c)=-1+\sum_{h=1}^H\frac{c_h}{d_i-b_h},\qquad 1\leq i\leq H.
\]
Their common zero set has exactly $H!$ connected components, each
a smooth manifold of dimension $H$. Every common zero has distinct
pole coordinates, and its residues are uniquely given by
\[
 c_h=-\frac{\prod_{i=1}^H(b_h-d_i)}
                  {\prod_{g\ne h}(b_h-b_g)}.
 \tag{1}
\]
\end{theorem}

\begin{proof}
Suppose first that a common zero has only $K<H$ distinct pole values.
Collect equal poles by adding their residues. Subtracting one and
clearing the $K$ distinct denominator factors gives a polynomial
in a scalar argument $d$ of degree exactly $K$, with leading
coefficient $-1$. Its $H$ distinct roots $d_i$ are impossible.
Thus all pole coordinates at a common zero are distinct.

For distinct $b_h$, consider the exact rational identity
\[
 -1+\sum_{h=1}^H\frac{c_h}{d-b_h}
 =-\frac{\prod_{i=1}^H(d-d_i)}{\prod_{h=1}^H(d-b_h)},
 \tag{2}
\]
with the residues from (1). Both sides tend to $-1$ as $d$ tends
to infinity. Their residues at each simple pole $b_h$ coincide.
Their difference is therefore a rational function with no finite
poles, tending to zero at infinity; it is identically zero.
Equivalently one may clear denominators: the difference numerator
has degree at most $H-1$ and vanishes at all $H$ distinct $b_h$.
Equation (2) proves that (1) is a common zero.

Uniqueness of the residues follows by the same root argument:
the difference of two interpolants has numerator degree at most
$H-1$ and vanishes at the $H$ distinct arguments $d_i$.
The common zero set is consequently the graph of (1) over
\[
 \{b\in(0,\beta)^H:b_h\ne b_g\text{ for }h\ne g\}.
\]
This domain has exactly $H!$ connected components, its convex
strict-order chambers. The graph projection is a homeomorphism,
as its inverse is the continuous map (1).
Finally the derivative of $(F_i)_i$ with respect to $(c_h)_h$
is the Cauchy matrix $(1/(d_i-b_h))_{i,h}$, which is nonsingular
by the same uniqueness argument. Thus all zero intersections
are regular, of codimension $H$ in dimension $2H$.
\end{proof}

\begin{corollary}[Whole-cube realization with logarithmically many row bits]
For $r\geq1$ and $S\geq1$, there is a family of $H=2^r$ actual
admissible affine-ratio heads on $n=r+S$ Boolean coordinates, with
$2HS$ free effective parameters, for which the zero equations on
the $HS$ inputs
\[
 (u,e_j),\qquad u\in\{0,1\}^r,\quad 1\leq j\leq S,
\]
have exactly $(H!)^S$ connected components. All denominators exceed
$3/4$ on the entire $n$-cube, throughout the parameter domain.
\end{corollary}

\begin{proof}
Number the row coordinates from $0$ to $r-1$ and set
\[
 D(u)=1-\sum_{i=0}^{r-1}2^{-i-4}u_i.
\]
Its $2^r$ values are distinct, by uniqueness of finite binary
expansions, and all exceed $7/8$. For each $h,j$ choose
$b_{hj}\in(0,1/(8S))$ and $c_{hj}\in\mathbb R$. Use
\[
 B_h(u,v)=D(u)-\sum_{j=1}^S b_{hj}v_j,
 \qquad A_h(u,v)=\sum_{j=1}^S c_{hj}v_j,
 \qquad F(u,v)=-1+\sum_h\frac{A_h(u,v)}{B_h(u,v)}.
\]
Every denominator slope is strictly negative and
$B_h>7/8-1/8=3/4$ on the whole cube. Hence these are admissible
affine-ratio denominators in the original model; their affine
numerators may have arbitrary real coefficients.
At $(u,e_j)$ the score equals
$-1+\sum_h c_{hj}/(D(u)-b_{hj})$. The equations for distinct
columns involve disjoint parameter pairs $(b_{\cdot j},c_{\cdot j})$.
Their common zero set is the Cartesian product of $S$ zero sets
from the theorem, so it has $(H!)^S$ connected components.
\end{proof}

For example, $S=r$ gives $n=2r$, $H=2^r\gg n$, and $p=2Hr$
free parameters. Since $\log(H!)=(1-o(1))H\log H$, the logarithm
of the component count is $(1-o(1))p\log H/2$. Thus a bound
$\exp(O(p))$ on all such regular zero-set intersections is false,
even in this regime with exponentially many heads and denominators
uniformly far from zero. Quotienting by a single simultaneous
permutation of all head labels leaves $(H!)^{S-1}$ component orbits,
so global head-label redundancy does not explain the growth.

This construction does not establish shattering of more than a
parameter-linear number of points. Symmetric extra row queries still
have the form (2), whose sign away from the prescribed zero rows is
independent of the pole ordering. The residue attached to the
$k$th smallest pole has sign $(-1)^{k+1}$, which supplies an
explicit ordering-dependent observable but not an independent
binary encoding of all orders. A general asymmetric decoder remains
unproved.

\section{Why denominator clearing cannot have bounded parameter degree}

\begin{theorem}[Necessary degree in the original effective coordinates]
Fix $H\geq 1$ and define, for $b_h>0$ and arbitrary real $a_h,c$,
\[
 F(a,b,c)=c+\sum_{h=1}^H\frac{a_h}{1+b_h}.
\]
If a real polynomial $Q(a,b,c)$ satisfies
\[
 \operatorname{sgn} Q(a,b,c)=\operatorname{sgn} F(a,b,c)
 \quad\hbox{whenever }b_h>0\hbox{ and }F(a,b,c)\ne0,
\]
then $Q$ has total degree at least $H+1$.
\end{theorem}

\begin{proof}
Write $z_h=1+b_h$ and
\[
 P(a,z,c)=c\prod_hz_h+\sum_h a_h\prod_{g\ne h}z_g.
\]
This polynomial is irreducible in
$\mathbb R[z_1,\ldots,z_H,a_1,\ldots,a_H,c]$.
Indeed, view it as a polynomial in $(a,c)$ over
$\mathbb R[z_1,\ldots,z_H]$. Its coefficients have greatest common
divisor one, because no $z_h$ divides the coefficient of $a_h$.
It is therefore primitive. Over the fraction field
$\mathbb R(z_1,\ldots,z_H)$ it is a nonconstant polynomial of degree
one in $(a,c)$, so it is irreducible there. Gauss's lemma gives the
claimed irreducibility over the polynomial ring.

For any fixed $a$ and $z_h>1$, let
$c_0=-\sum_h a_h/z_h$. The score changes sign at $c=c_0$.
Continuity of $Q$ and the required signs on both sides force
$Q(a,z-1,c_0)=0$. Substitute this rational value of $c$ into $Q$.
After clearing powers of $\prod_hz_h$, the result is a polynomial
in $(a,z)$ vanishing on the open set $z_h>1$, so it is identically
zero. Polynomial division in $c$ over the fraction field
$\mathbb R(a,z)$ now shows that $c+\sum_h a_h/z_h$ divides $Q$
in that field. Equivalently, $P$ divides $Q$ after localization by
the nonzero polynomials in $(a,z)$. Since $P$ is irreducible and
depends on $c$, it divides none of those inverted polynomials;
primality in the polynomial ring therefore gives $P\mid Q$ already
before localization. Translation $b=z-1$ preserves total degree.
Finally $P$ has total degree $H+1$, as seen from its term
$c\prod_hz_h$, and $Q$ is nonzero because it has both required
strict signs. Thus $\deg Q\geq H+1$.
\end{proof}

These scores are evaluations of sums of actual admissible affine
ratios. For example, on the one-bit cube take numerator $a_h$ and
denominator $1+b_h x$, whose slope is strictly positive. The theorem
does not constrain a nonlinear global reparameterization. It rules
out replacing a sum of independently parameterized fractions by a
sign polynomial of degree bounded independently of $H$ in those same
effective coordinates, even when every denominator is positive.

\section{Fourier and finite-difference rank do not count heads}

\begin{theorem}[A single reciprocal has full transformed flattening rank]
Fix $r\geq1$ and $k\geq0$. For $u,v\in\{0,1\}^r$ put
$t(u)=\sum_{i=0}^{r-1}2^i u_i$. Fix $a_1,\ldots,a_k>0$ and consider
the single admissible affine ratio
\[
 R(u,v,z)=\frac1{1+t(u)+t(v)+\sum_{j=1}^k a_jz_j},
 \qquad z\in\{0,1\}^k.
\]
Let $M$ be the $2^r$ by $2^r$ matrix obtained by taking its mixed
forward difference in all $k$ coordinates of $z$ at zero:
\[
 M_{u,v}=\sum_{z\in\{0,1\}^k}(-1)^{k-|z|}R(u,v,z).
\]
Then $(-1)^kM$ is strictly positive definite and has rank $2^r$.
Applying the Boolean Walsh transform, the Boolean M\"obius transform,
or any invertible linear transform independently to the $u$ and $v$
axes preserves this full rank.
\end{theorem}

\begin{proof}
Every denominator slope is strictly positive and its constant term
is one, so the ratio is admissible on the entire cube. The elementary
Laplace identity $q^{-1}=\int_0^\infty e^{-qs}\,ds$, valid for $q>0$,
gives
\[
 (-1)^k M_{u,v}
 =\int_0^\infty e^{-s}e^{-t(u)s}e^{-t(v)s}
          \prod_{j=1}^k(1-e^{-a_js})\,ds.
\]
For $k=0$ the product is one. For every $k$ the density multiplying
$e^{-t(u)s}e^{-t(v)s}$ is strictly positive for $s>0$.
Thus for any nonzero vector $(q_u)$,
\[
 (-1)^k q^{\mathsf T}Mq
 =\int_0^\infty e^{-s}\prod_{j=1}^k(1-e^{-a_js})
       \left(\sum_u q_u e^{-t(u)s}\right)^2\,ds>0.
\]
To justify strictness, the integers $t(u)$ are all distinct, so the
nonzero polynomial $\sum_u q_u y^{t(u)}$ cannot vanish on all
$y\in(0,1)$. Substituting $y=e^{-s}$ shows that the squared factor is
positive on some open interval. The integral is finite, since it is
bounded by a constant times $e^{-s}$. Strict positive definiteness
implies full rank. Finally, independent invertible transforms replace
$M$ by $UMV^{\mathsf T}$ for invertible matrices $U,V$, which preserves
rank. Both named Boolean transforms are invertible.
\end{proof}

This counterexample permits arbitrary signed affine numerators in the
ambient model but needs only a constant positive numerator. Ordinary
tensor rank bounds obtained from these flattenings therefore cannot
assign $O(n)$ rank to each head. The conclusion remains true after
arbitrarily many mixed Boolean differences in additional coordinates.
It does not rule out an invariant using positivity or a joint
factorization shared by several transforms.

\begin{theorem}[Fourier threshold density is not a polynomial head invariant]
For odd $m\geq3$, let $f_m$ be the conjunction of majority on two
disjoint blocks of $m$ bits. In the unchanged attention model,
\[
 H^*(f_m)\leq 2m+2.
\]
Nevertheless its Fourier threshold density is $m^{\Omega(\log m)}$:
every strict polynomial sign representation in the Walsh characters
requires that many nonzero coefficients. Consequently no fixed
polynomial in $nH^*(f)$ bounds Fourier threshold density for all
Boolean functions on $n$ bits.
\end{theorem}

\begin{proof}
Write $\sigma_k\in\{-1,1\}$ for majority at Hamming weight $k$.
Fix distinct positive
numbers $\lambda_0,\ldots,\lambda_m$. The square Cauchy matrix
$C_{k,j}=1/(k+\lambda_j)$, $0\leq k,j\leq m$, is invertible.
Indeed, a vector in its kernel gives a rational function
$\sum_j c_j/(t+\lambda_j)$ vanishing at $m+1$ distinct points.
After multiplication by $\prod_j(t+\lambda_j)$ its numerator has
degree at most $m$, so it is zero. Taking the residue at each pole
then gives $c_j=0$ for all $j$.

Choose the unique vector $c$ with $Cc=\sigma$ and put
\[
 U(t)=\sum_{j=0}^m\frac{c_j}{t+\lambda_j}.
\]
On the Boolean cube the score
$U(|x|)+U(|y|)-1$ equals $1$ when both majorities are positive,
$-1$ when exactly one is positive, and $-3$ otherwise. It has strict
margin at least one. The denominators in its $2m+2$ displayed terms
have positive slopes in the active block and zero slopes in the
other block. Replace them, respectively, by
\[
 |x|+\lambda_j+\varepsilon|y|,
 \qquad |y|+\lambda_j+\varepsilon|x|.
\]
For every $\varepsilon>0$ all slopes are strictly positive, and each
denominator is at least $\lambda_j>0$ on the entire cube. Each ratio
therefore compiles to one actual admissible head. The perturbed
score converges to the preceding score at every cube point as
$\varepsilon\downarrow0$. Since the cube is finite, convergence is
uniform there; choose $\varepsilon>0$ with error less than $1/2$.
Every required strict sign is preserved, proving the head bound.

The density lower bound is Sherstov's Theorem 6.3 in
\emph{The Intersection of Two Halfspaces Has High Threshold Degree},
\url{https://web.cs.ucla.edu/~sherstov/pdf/hshs.pdf}. Its density is
defined using monomials on $\{-1,1\}^{2m}$, exactly the Walsh
characters after the coordinate change from $\{0,1\}^{2m}$.
The theorem states
$\operatorname{dns}(\operatorname{MAJ}_m\wedge
\operatorname{MAJ}_m)=m^{\Omega(\log m)}$. This outgrows every fixed
polynomial in $(2m)(2m+2)$, proving the last assertion.
\end{proof}

The Cauchy compilation is a complete written construction. The density
lower bound in this theorem is an application of the cited published
result, not a new proof or a Lean formalization. Its role here is to
reject a concrete proposed invariant; a restricted-class upper bound
does not resolve the worst-case asymptotic question.

\section{Equal moments limit exterior Stieltjes readout}

The next theorem addresses a proposed way to keep the first several
terms of a compact-pole expansion fixed while encoding information in
higher moments. It permits arbitrary correlated choices of a pole
multiset inside one queried slot block. It does not require independent
pair choices or a fixed finite library of poles.

\begin{proposition}[A zero-count criterion from prefix Wronskians]
Let $f_0,\ldots,f_s$ be smooth real functions on an open interval $I$.
Suppose that every prefix Wronskian
$W(f_0,\ldots,f_j)$, $0\leq j\leq s$, is everywhere nonzero on $I$.
Then each nonzero linear combination of the $f_j$ has at most $s$
distinct zeros in $I$.
\end{proposition}
\begin{proof}
We use induction on $s$. For $s=0$ the assertion is exactly the
nonvanishing of $f_0$. For the inductive step divide a proposed
combination by the nowhere-zero function $f_0$ and differentiate.
The derivative is a linear combination of
\[
 g_j=(f_{j+1}/f_0)',\qquad 0\leq j<s.
\]
The elementary Wronskian identities for a common factor and a constant
first column give
\[
 W(g_0,\ldots,g_j)
 =\frac{W(f_0,\ldots,f_{j+1})}{f_0^{j+2}}\ne0.
\]
For completeness, multiplying all columns of a Wronskian of length
$m$ by a smooth function $a$ multiplies its determinant by $a^m$.
This follows from the product rule by expressing the transformed
derivative rows as a triangular linear combination of the old rows,
with diagonal entries $a$. Apply this with $a=f_0$ to
$1,f_1/f_0,\ldots,f_{j+1}/f_0$, and expand along the initial
constant column, to obtain the displayed identity.

If all coefficients other than that of $f_0$ vanish, the original
combination has no zeros. Otherwise the derivative combination is
nonzero, because a nonzero full Wronskian implies linear independence.
By induction this derivative has at most $s-1$ distinct zeros. If the
original combination had $s+1$ distinct zeros, so would its quotient
by $f_0$, and Rolle's theorem would give at least $s$ distinct derivative
zeros, a contradiction. The argument also excludes an infinite set of
zeros, since it would contain $s+1$ distinct zeros.
\end{proof}

\begin{theorem}[Equal-moment Stieltjes capacity]
Fix integers $k\geq1$ and $0\leq\ell\leq k-1$, and a real number
$\beta$. Let $\mathcal S$ be any family of multisets of exactly $k$
real poles, all at most $\beta$. Suppose all members of $\mathcal S$
have the same power sums through degree $\ell$:
\[
 \sum_{\lambda\in S}\lambda^j
 =\sum_{\lambda\in T}\lambda^j
 \quad(S,T\in\mathcal S,\ 1\leq j\leq\ell).
\]
Multiplicity is included in every sum. Put
$F_S(D)=\sum_{\lambda\in S}(D-\lambda)^{-1}$ for $D>\beta$.
For fixed query points $D_1,\ldots,D_r>\beta$ and fixed arbitrary
thresholds $\theta_1,\ldots,\theta_r$, suppose that
\[
 \bigl(\operatorname{sgn}(F_S(D_i)-\theta_i)\bigr)_{i=1}^r,
 \qquad S\in\mathcal S,
\]
realizes every element of $\{-1,1\}^r$ with strict signs. Then
\[
 r\leq k-\ell.
\]
In particular, without any moment conditions the bound is $r\leq k$.
If all first $k-1$ power sums are fixed, at most one query can be
shattered.
\end{theorem}
\begin{proof}
Two equal query locations cannot be shattered: whichever of their two
thresholds is smaller forbids one of the two opposite strict sign
patterns. Thus the query locations are distinct, and we may order them
as $D_1<\cdots<D_r$.

For $S,T\in\mathcal S$ define the monic degree-$k$ polynomials
\[
 p(D)=\prod_{\lambda\in S}(D-\lambda),\qquad
 q(D)=\prod_{\lambda\in T}(D-\lambda).
\]
Newton's identities show that matching power sums through degree
$\ell$ implies matching elementary symmetric sums through degree
$\ell$. Hence $u=q-p$ has degree at most
$s=k-\ell-1$. Both $p$ and $q$ are strictly positive on
$(\beta,\infty)$, and logarithmic differentiation gives
\[
 F_S-F_T=\frac{p'}p-\frac{q'}q
 =-\frac{(u/p)'}{1+u/p}.
\]
The denominator $1+u/p=q/p$ is strictly positive. We claim that a
nonzero numerator $(u/p)'$ has at most $s$ distinct zeros on this
interval.

For $0\leq j\leq s$ put $f_j=(D^j/p)'$. Every prefix Wronskian is
\[
 \begin{aligned}
 W(f_0,\ldots,f_j)
 &=W(1,1/p,D/p,\ldots,D^j/p)\\
 &=p^{-j-2}W(p,1,D,\ldots,D^j)\\
 &=(-1)^{j+1}\left(\prod_{a=0}^j a!\right)
       \frac{p^{(j+1)}}{p^{j+2}}.
 \end{aligned}
\]
The first equality expands along a constant first column. The second
is the common-factor identity proved above. For the last equality,
move the first column to the last position, then expand the resulting
Wronskian of $1,D,\ldots,D^j,p$ along its final derivative row.
Here $p^{(j+1)}$ denotes a derivative. It is strictly positive for
$D>\beta$: differentiating a product of $k$ strictly positive linear
factors $j+1\leq s+1\leq k$ times gives a sum of positive products,
with positive integer coefficients. Thus all these prefix Wronskians
are nonzero, including in the presence of repeated poles. The
preceding proposition proves the claim for every polynomial $u$ of
degree at most $s$.

Now choose $S$ giving the alternating strict sign pattern
$((-1)^i)_{i=1}^r$ relative to the thresholds, and $T$ giving its
negative. Their difference $F_S-F_T$ has the same alternating signs
at the $D_i$. It is nonzero and has at least $r-1$ distinct zeros,
one in each intervening open interval. The claim gives
$r-1\leq s=k-\ell-1$, as required.
\end{proof}

The theorem concerns unweighted Stieltjes sums on a single slot block.
It allows the poles to move continuously and correlatedly, but it does
not sum over several overlapping slot blocks. Applying the theorem
separately to disjoint blocks with a total of $H$ poles bounds their
total number of independently readable strict bits by $H$. Applying
it separately to nested blocks can instead give an $H\log H$ sum,
because the same poles are counted at multiple levels. Whether a
hereditary family of correlated equal-moment partitions can make those
bits simultaneously readable remains a separate unresolved question.
The present theorem supplies no asymptotic bound on unrestricted head
complexity.

\section{Primary-source transfer checks}

Pardo and Sebasti\'an, \emph{Erzeugungsgrad, VC-Dimension and Neural
Networks with rational activation function},
\url{https://arxiv.org/html/2504.11345v1}, Corollary 5.13, studies
classifiers given by distinguished Zariski open sets
$D(f/g)=\{x:f(x)g(x)\ne0\}$ over algebraically closed fields.
Strict positivity on the real cube is a different classifier.
The cited result therefore does not supply the needed sign-pattern
bound; its degree factors also remain in the logarithm.

Bartlett and Williamson, \emph{The VC Dimension and Pseudodimension
of Two-Layer Neural Networks with Discrete Inputs}, Neural Computation
8(3) (1996), 625--628,
\url{https://doi.org/10.1162/neco.1996.8.3.625}, states a
$W\log(WD)$ bound for the indicated sigmoid and radial-basis networks.
This is consistent with the ordinary denominator-clearing degree
growth and does not give a parameter-linear bound for attention.
The frequently quoted fan-in refinement in Anthony and Bartlett's
book, Theorem 8.11, was not used here: its printed proof asserts a
degree bound independent of the number of summed rational units,
and that step requires an additional argument. The preceding theorem
shows why simply copying such a degree assertion into the affine-ratio
model is invalid. No claim that the book's theorem itself is false
is needed for this transfer audit.

\paragraph{Remaining lower-bound gap.}
None of these arguments controls the sign patterns of unrestricted
sums with target-dependent denominators. A successful opposite-direction
resolution still needs a full-cube obstruction forcing at least
$c2^n/n$ heads for some Boolean table, uniformly in all admissible
coefficients. Fixed-denominator dimension bounds, arbitrary-sample
parameter bounds, and symmetric-sheet restrictions do not supply it.
\end{document}
